\section{Projects}
% \subsection{Image Generation using Multi-marginal Optimal Transport}{University of Toronto}{2022 - 2023}
% \begin{itemize}[leftmargin=*]
%     \item Developed an advanced image generation framework leveraging multi-marginal optimal transport theories, utilizing PyTorch. This project emphasizes my capability to translate complex theoretical models into practical computational algorithms.
%     \item Analyzed the CelebA dataset, comprising over 200,000 images, applying sophisticated numerical optimization techniques to improve generative model performance.
% \end{itemize}
\subsection{Population Model Using SDE and Game Theory}{University of Toronto}{2021 - 2023}
\begin{itemize}[leftmargin=*]
    \item PhD research project on the generative modeling of scRNA-seq data using stochastic differential equations and game theory.
    \item Published findings in PLOS Computational Biology; developed novel algorithms for high-dimensional population dynamics modeling. (\href{https://github.com/Diebrate/population_model}{manuscript} and \href{https://journals.plos.org/ploscompbiol/article?id=10.1371/journal.pcbi.1012015}{code})
\end{itemize}

\subsection{Anomaly Detection Using Optimal Transport}{University of Toronto}{2021 - 2023}
\begin{itemize}[leftmargin=*]
    \item PhD research project on developing statistical methods for anomaly detection in time series data, utilizing principles of optimal transport closely related to WGAN.
    \item Designed specialized numerical algorithms that improved detection accuracy and computational efficiency. (\href{https://arxiv.org/abs/2501.03501}{manuscript} and \href{https://github.com/Diebrate/scRNA_study}{code})
\end{itemize}
